\section{梯度消失/爆炸与 LSTM}

\subsection{梯度消失问题的根源}

反向传播中，隐藏状态之间的梯度为：

$$\frac{\partial h_t}{\partial h_{t-1}} = \text{diag}(\tanh'(\cdot)) \cdot W_{hh}$$

这个乘积在时间步之间\textbf{反复累乘}，导致：

\begin{figure}[H]
\centering
\includegraphics[width=0.95\textwidth]{slides-images/slide-048.jpg}
\caption{梯度消失与爆炸的两个来源：$\tanh$ 导数的最大值为 1（几乎总是小于 1），以及 $W_{hh}$ 矩阵的最大奇异值可能大于或小于 1\protect\footnotemark}
\end{figure}
\footnotetext{Slides 第49页。}

\begin{importantbox}{梯度消失/爆炸的数学分析}
\begin{enumerate}
    \item \textbf{$\tanh$ 导数}：$\tanh'(x) \in [0, 1]$，只有在 $x = 0$ 时取最大值 1。多次乘以小于 1 的值 $\rightarrow$ 指数衰减（梯度消失）。
    \item \textbf{$W_{hh}$ 的奇异值}：
    \begin{itemize}
        \item 最大奇异值 $\sigma_{\max} > 1$：梯度爆炸
        \item 最大奇异值 $\sigma_{\max} < 1$：梯度消失
    \end{itemize}
\end{enumerate}
梯度爆炸可以通过\textbf{梯度裁剪}（gradient clipping）解决：当梯度范数超过阈值时进行缩放。但梯度消失更加棘手——这正是 LSTM 要解决的问题。
\end{importantbox}

\subsection{LSTM 的核心思想}

LSTM（Long Short-Term Memory）通过引入\textbf{门控机制}和\textbf{细胞状态}（cell state）来缓解梯度消失问题。

\begin{figure}[H]
\centering
\includegraphics[width=0.95\textwidth]{slides-images/slide-049.jpg}
\caption{LSTM 结构图：四个门（input gate, forget gate, output gate, gate gate）控制信息的写入、遗忘和读出。顶部的细胞状态通道（cell state highway）提供了无激活函数干扰的信息传递路径\protect\footnotemark}
\end{figure}
\footnotetext{Slides 第50页。}

LSTM 的四个门：
\begin{enumerate}
    \item \textbf{Gate Gate（候选值）}：计算要写入的候选信息 $\tilde{c}_t = \tanh(W_g [h_{t-1}, x_t])$
    \item \textbf{Input Gate}：决定是否写入 $i_t = \sigma(W_i [h_{t-1}, x_t])$
    \item \textbf{Forget Gate}：决定遗忘多少历史信息 $f_t = \sigma(W_f [h_{t-1}, x_t])$
    \item \textbf{Output Gate}：决定输出多少 $o_t = \sigma(W_o [h_{t-1}, x_t])$
\end{enumerate}

细胞状态更新：$c_t = f_t \odot c_{t-1} + i_t \odot \tilde{c}_t$

隐藏状态输出：$h_t = o_t \odot \tanh(c_t)$

\begin{importantbox}{LSTM 解决梯度消失的关键：Cell State Highway}
细胞状态 $c_t$ 的更新路径（图中顶部通道）\textbf{不经过 $\tanh$ 激活函数}，只涉及逐元素乘法（forget gate）和加法（input gate）。这意味着：
\begin{itemize}
    \item 只要 forget gate 不为零，梯度就能相对无损地沿时间反向传播
    \item 类似于 ResNet 中的跳跃连接——提供了梯度的``高速公路''
    \item LSTM 并不保证解决梯度消失，但使得模型\textbf{更容易}学习长距离依赖
\end{itemize}
\end{importantbox}

\subsection{LSTM 与 ResNet 的联系}

\begin{figure}[H]
\centering
\includegraphics[width=0.95\textwidth]{slides-images/slide-051.jpg}
\caption{LSTM 的 cell state highway 与 ResNet 的跳跃连接在概念上高度相似：都通过``绕过''非线性层来改善梯度流动\protect\footnotemark}
\end{figure}
\footnotetext{Slides 第52页。}

\begin{knowledgebox}{跳跃连接的普适原理}
ResNet 和 LSTM 解决的问题在本质上是相同的：
\begin{itemize}
    \item \textbf{ResNet}：随着网络\textbf{深度}增加，梯度消失使得深层难以训练 $\rightarrow$ 引入跳跃连接
    \item \textbf{LSTM}：随着序列\textbf{长度}增加，梯度消失使得远距离信息难以利用 $\rightarrow$ 引入 cell state highway
\end{itemize}
核心思想是一致的：提供一条不经过非线性变换的``直达路径''，让信息和梯度能够更自由地流动。
\end{knowledgebox}

\subsection{GRU 简介}

GRU（Gated Recurrent Unit）是 LSTM 的一个简化变体，将 forget gate 和 input gate 合并为一个 update gate，减少了参数量，同时保持了类似的性能。

\subsection{本章小结}

梯度消失是 Vanilla RNN 难以处理长序列的根本原因。LSTM 通过引入门控机制和无激活函数的 cell state 通道来缓解这一问题。这与 ResNet 中的跳跃连接异曲同工——都通过提供梯度``高速公路''来改善深层/长距离的训练。

%% ========== 状态空间模型 ==========

